\section{Planteamiento del problema}

La visualización de campos de flujo, mediante técnicas como Streamlines, Streaklines y Pathlines, es fundamental en diversas áreas de la ciencia e ingeniería, incluyendo la dinámica de fluidos computacional (CFD), la meteorología, y el modelado de procesos físicos \cite{camp_streamline_2011}. Estos métodos permiten representar de forma intuitiva el comportamiento de partículas en movimiento dentro de un campo vectorial, lo cual facilita el análisis y la interpretación de fenómenos complejos.



Sin embargo, la simulación precisa y eficiente de estas trayectorias implica un alto costo computacional, especialmente cuando se trabaja con volúmenes de datos grandes o se requieren resoluciones espaciales y temporales finas \cite{pugmire_scalable_2009}. En este contexto, la programación tradicional en C/C++, aunque eficiente, puede resultar insuficiente para explotar de manera óptima las capacidades de cómputo de los sistemas modernos, en especial aquellos que cuentan con aceleradores como GPUs.



En respuesta a esta necesidad, han surgido enfoques como la programación heterogénea, que busca distribuir la carga computacional entre distintos dispositivos de procesamiento, permitiendo una ejecución más rápida y eficiente. En particular, el estándar SYCL, basado en C/C++, ofrece una solución portable y flexible para desarrollar aplicaciones que aprovechen arquitecturas heterogéneas sin perder la expresividad y control del lenguaje de bajo nivel \cite{keryell_sycl_2019}.


A pesar de sus beneficios teóricos, existe una brecha de conocimiento práctico respecto al impacto real del uso de SYCL en problemas computacionalmente intensivos como el cálculo de líneas de flujo. No se cuenta con suficiente evidencia comparativa que cuantifique su desempeño frente a enfoques tradicionales en términos de tiempo de ejecución, consumo de memoria y escalabilidad, con respecto a esta técnica de visualización en particular. Para cuantificar el desempeño se realizan estudios experimentales usando métricas como tiempo de ejecución, uso de memoria y escalabilidad \cite{camp_gpu_2013}, relacionadas con características del problema como: tamaño del dataset, tamaño del conjunto de semillas, distribución del conjunto de semillas y complejidad del campo vectorial \cite{camp_streamline_2011}.
